% !TEX root = ../Tesis.tex
\chapter{Conclusión}

La presente investigación evaluó el impacto de la contaminación atmosférica por 
procesos de transformación y combustión de hidrocarburos a través del 
fraccionamiento isotópico de $\delta^{18}$O y $\delta$D en el vapor de agua 
atmosférico de Salamanca, Guanajuato. A partir de la integración del análisis 
isotópico, la caracterización de contaminantes criterio y la modelación 
lagrangiana de dispersión, se concluye lo siguiente:

\begin{itemize}
    \item  Se acepta la hipótesis de que 
    el vapor de agua emitido durante la transformación y combustión de 
    hidrocarburos altera la composición isotópica del vapor atmosférico en 
    Salamanca. El modelo binario de mezcla estimó una contribución de vapor 
    derivado de la combustión (VDC) de entre 18.9 y 21.5\,\%, y los valores 
    negativos de \textit{d-excess} (hasta $-38.85$\,\textperthousand) 
    confirmaron la presencia de una componente isotópica distinta a la del 
    ciclo hidrológico natural. Por ello, se valida el uso del deuterio como 
    trazador de la intensidad de la contaminación antropogénica.

    \item Se identificó que las concentraciones de 
    contaminantes criterio están dominadas por la meteorología local. Las 
    condiciones estables de febrero favorecieron la acumulación de emisiones 
    primarias (mayor NO$_2$ y PM$_{10}$), mientras que la mayor temperatura y 
    humedad de mayo propiciaron una atmósfera más mezclada. El análisis de 
    correlación de Spearman y los ciclos diurnos permitieron atribuir el NO$_2$ 
    al tráfico vehicular y el SO$_2$ a fuentes industriales fijas, evidenciando 
    una influencia mixta en la ciudad.

    \item Las 
    simulaciones de FLEXPART con CO como trazador indicaron una influencia 
    predominantemente local de las emisiones, con una mayor sensibilidad en las 
    celdas contiguas a cada receptor. La similitud de los \textit{footprints} 
    entre estaciones sugiere un origen común del vapor de fondo, mientras que 
    la reducción del área de influencia nocturna en mayo indicó una menor 
    probabilidad de que las emisiones de las industrias del sureste alcanzaran 
    los receptores. Se identificaron como fuentes potenciales la refinería, la 
    termoeléctrica y empresas de movilización de combustible, ubicadas al 
    noreste de la ciudad.

    \item Se determinó que los valores de $\delta^{18}$O y 
    $\delta$D del vapor atmosférico presentan una variación temporal 
    consistente con la influencia de emisiones locales. Las muestras de febrero 
    se asociaron con valores de \textit{d-excess} más negativos ($\approx 
    -30$\,\textperthousand) y mayores concentraciones de PM$_{2.5}$, PM$_{10}$, 
    CO y NO$_2$, mientras que las de mayo mostraron valores menos negativos 
    ($\approx -24$ a $-31$\,\textperthousand), consistente con una menor 
    contribución relativa de VDC.

    \item Se evaluó que la variación del 
    fraccionamiento isotópico está modulada por la carga de contaminantes 
    primarios y las condiciones de estabilidad atmosférica. La muestra de 
    febrero con el \textit{d-excess} más alto ($\approx -30$\,\textperthousand) 
    coincidió con los mayores valores de PM$_{2.5}$, PM$_{10}$, CO y NO$_2$, y 
    el menor de O$_3$, lo que sugiere que las noches de mayor acumulación de 
    emisiones locales son las que registran la señal isotópica de combustión 
    más clara.
\end{itemize}

Los resultados deben interpretarse 
considerando que la estimación de la contribución del VDC se realizó con una 
firma isotópica de referencia de origen teórico. Por ello, resulta fundamental 
determinar experimentalmente la composición isotópica del vapor generado por 
la combustión bajo condiciones representativas de las fuentes y condiciones 
ambientales locales, con el propósito de reducir la incertidumbre asociada al 
modelo de mezcla. Asimismo, es necesario ampliar el número de muestras y los 
periodos de observación para evaluar la variabilidad temporal y espacial de la 
composición isotópica del vapor atmosférico.

Finalmente, esta investigación constituye una base para incorporar el análisis 
isotópico del vapor de agua atmosférico al estudio de la contaminación urbana e 
industrial en México, abriendo nuevas perspectivas para el desarrollo de 
herramientas complementarias de identificación y evaluación de la influencia 
antropogénica sobre la calidad del aire y la composición atmosférica a escala 
local.